\scanpage{70}
où $\mathcal F_X^{\mathrm{an}*}$ se \deleted{\unsure{construit}} définit à partir de $X^{\mathrm{an}}$ \deleted{par}\inserted{de} façon analogue que dans le cas algébrique :
\[
\mathcal F_{X^{\mathrm{an}}}^{\deleted{\unsure{\mathrm{an}}}\mathrm{an}*}
=\mathcal F_{X^{\mathrm{an}}}^*
=\deleted{\unsure{\mathcal F}}\mathcal F_{X^{\mathrm{an}}}^{\nu}
\]
\[
\mathcal F_{X^{\mathrm{an}}}^{\nu}:\quad U\longmapsto\Gamma(U^{\nu+1}_{/U},\mathcal O_{U^{\nu+1}_{/U}})
\]
(où $U$ parcourt les ouverts de $X^{\mathrm{an}}$).

\uline{Conjecture} Pour tout espace analytique $\mathfrak X$, le complexe $\mathcal F_{\mathfrak X}^*$ est une \uline{résolution} du faisceau $\mathbf C_{\mathfrak X}$ (via l'augmentation triviale).
\margin{[Alors pour $X$ prop, on a $H^*(\mathfrak X_{\mathrm{str}},\mathcal O_{\mathfrak X_{\mathrm{str}}})=H^*(\mathfrak X,\mathbf C)$]}
\deleted{\uline{NB} \unsure{Cela} \illegible{9}}
Cela semble plausible, en regardant les ``fibres formelles'' des faisceaux...

\uline{Th} Soit $X$ propre sur $\mathbf C$. Si la conjecture plus haut est vraie pour $X^{\mathrm{an}}$, alors on a un isom can
\[
H^*(X_{\mathrm{strat}},\mathcal O_{X_{\mathrm{strat}}})\simeq H^*(X^{\mathrm{an}},\mathbf C)
\]
\mspage{69}
